\section{Robustness and Extensions}
\label{sec:robustness}

We subject our findings to extensive robustness checks organized in six
groups. Section~\ref{sec:robustness}.1 varies the FL definition
threshold, win-rate cutoff, and temporal window.
Section~\ref{sec:robustness}.2 checks sample restrictions and
specification variants (unrestricted sample, homogeneous subsamples,
tighter controls, matching, clustering).
Section~\ref{sec:robustness}.3 runs sensitivity analysis on unobservable
confounders. Section~\ref{sec:robustness}.4 reports the staggered DiD
(complementary, underpowered). Section~\ref{sec:robustness}.5 computes
welfare bounds under OLS, IV, and cross-fit estimates.
Section~\ref{sec:robustness}.6 examines procurement oversight and
compliance channels.

\subsection{FL Definition Robustness}

\paragraph{IQR threshold sensitivity.}
We vary the IQR multiplier from 1.0 to 3.0, which changes both the
threshold and the number of FL firms (from 3,442 at 1.0$\times$ to 1,456
at 3.0$\times$). The price coefficient increases monotonically with the % REVISION-V5: dual interpretation per referee minor 10
multiplier. This pattern is consistent with a stricter threshold
selecting more extreme cover bidders, but also with progressively more
extreme firms participating in more unusual (higher-priced) markets for
non-collusive reasons. We cannot distinguish between these
interpretations from the threshold sensitivity exercise alone.
Figure~\ref{fig:threshold_stability} and Table~\ref{tab:threshold} in
Appendix~\ref{app:robustness} report details.

\paragraph{Low-win-rate FL variants.}
Instead of requiring a zero win rate, we define FL using win rate
$< 1\%$ and $< 2\%$. These variants capture firms that win trivially
(1--2 wins in hundreds of participations) but otherwise behave as cover
bidders. Table~\ref{tab:fl_lowwinrate} in Appendix~\ref{app:robustness}
shows that coefficients are stable across these definitions.

\paragraph{Cross-fitting.}
We define FL firms using only odd years (2009, 2011, \ldots, 2019) and
estimate regressions on even years, and vice versa. This breaks any
mechanical link between FL classification and outcomes.
Table~\ref{tab:fl_crossfit} in Appendix~\ref{app:robustness} shows that
the odd-year-trained fold yields a significant coefficient of 0.053
($p < 0.05$); the even-year-trained fold yields 0.019, which is
positive but not individually significant, likely reflecting lower
statistical power from fewer FL firms identified in even years (1,885
vs.\ 2,153). The cross-fold average of 0.036 is attenuated relative to
the full-sample OLS (0.064) but remains positive, consistent with the
FL screen capturing a genuine signal.

\paragraph{Temporal FL definition.}
Instead of classifying FL firms over the full 2009--2019 period, we use
rolling 3-year windows: a firm is FL in year $t$ if it is an
always-loser with above-threshold participation in years $[t-2, t]$.
Table~\ref{tab:fl_temporal} in Appendix~\ref{app:robustness} shows
that time-varying FL status produces a similar coefficient of 0.070
($p < 0.01$, $N = 1{,}255{,}998$), close to the baseline 0.064.

\subsection{Sample and Specification Robustness}

\paragraph{Unrestricted sample.}
Our baseline restricts to item types with at least one FL-present
tender. Removing this restriction yields a coefficient of 0.063 ($N =
2{,}875{,}653$), virtually identical to the baseline 0.064
(Table~\ref{tab:unrestricted_sample} in Appendix~\ref{app:robustness}),
suggesting that restriction (iii) does not bias the results.

\paragraph{Homogeneous sub-samples.}
Identifying \emph{ex ante} homogeneous goods from item descriptions is
infeasible given BEC's coding structure. Instead, we use a data-driven
approach: for each item group (first two digits of the item code), we
compute the within-group coefficient of variation (CV) of log negotiated
prices. Groups with low CV contain items with similar prices across
tenders---a proxy for product homogeneity. We split item groups at the
median CV and run FL regressions on each subsample.
Table~\ref{tab:homogeneous_cv} in Appendix~\ref{app:robustness} reports
the results. If the FL effect is driven by quality differences rather
than cover bidding, we would expect it to vanish in low-CV (homogeneous)
groups where quality-based selection stories are implausible. Item fixed
effects already absorb much cross-item heterogeneity, so the CV split
provides an additional dimension of variation.

\paragraph{Tighter controls.}
We add controls beyond item/year/PBU FE.
Table~\ref{tab:tighter_controls} in Appendix~\ref{app:robustness}
reports four specifications. Adding the number of genuine bidders
increases the coefficient to 0.094 (column~2), consistent with FL
presence being positively selected into competitive tenders.
Item $\times$ year FE (column~4) yield 0.074, close to the baseline.
Column~(3) adds log reference price, which absorbs most cross-tender
price variation ($R^2$ rises from 0.886 to 0.973); the FL coefficient
becomes $-0.031$. This sign reversal is consistent with a bad-control
problem: if cover bidding inflates reference prices---e.g., through
prior-tender manipulation---then conditioning on reference price removes
the treatment effect and may induce collider bias. We therefore treat
column~(3) as a mechanical decomposition of the price level rather than
a robustness check.

\paragraph{Matching estimators.}
Coarsened exact matching (CEM) on year, procedure type, item group, and
PBU size quartile yields a coefficient of 0.077 ($N = 969{,}751$).
Inverse probability weighting (IPW) yields 0.055 ($N = 830{,}194$).
Both bracket the OLS estimate of 0.064 from the general+PBU
specification.

\paragraph{Alternative clustering.}
Standard errors clustered at the item level (baseline: 0.022), PBU level
(0.014), and two-way item + PBU level (0.024) all leave inference
unchanged, with the coefficient significant at the 1\% level under all
clustering schemes.

\subsection{Benchmarking Against Bid-Level Screens} % REVISION-V5: new subsection per referee M8

To assess whether the FL screen adds value relative to existing % REVISION-V5: horse-race regression per referee M8/condition 2
bid-level screens, we compute an Imhof-style indicator based on the
within-item-group coefficient of variation (CV) of bid prices
\citep{imhof2019detecting}. In a horse-race regression controlling for
both FL presence and the Imhof flag, the FL coefficient remains
significant and increases slightly (0.084, SE = 0.022) relative to the
baseline (0.064), while the Imhof flag captures a different dimension of
tender characteristics ($-0.213$, SE = 0.006). The correlation between
the two screens is low (0.060), confirming that FL and bid-level screens
capture different signals: FL identifies suspicious \emph{firms}, while
CV-based screens identify suspicious \emph{tender patterns}. Details are
in Appendix~\ref{app:heterogeneity}
(Table~\ref{tab:screen_horserace}).

\subsection{Sensitivity to Unobservables}

\cites{cinelli2020making} sensitivity framework quantifies how strong
omitted variable bias would need to be to explain away the FL
coefficient. The robustness value $RV_{q=1} = 17.5\%$ means that an
unobserved confounder would need to explain at least 17.5\% of the
residual variation in both FL presence and log prices to reduce the
coefficient to zero---a stringent requirement given the rich fixed
effects structure.

\subsection{Staggered Difference-in-Differences}

A staggered DiD using \citet{callaway2021difference} is underpowered % REVISION-V5: DiD moved to appendix per referee M6
for effects of the magnitude we estimate via OLS (MDE~$= 0.076 >$~OLS
of 0.064) and is reported in Appendix~\ref{app:did} for completeness.

\subsection{Welfare Analysis}
\label{sec:welfare}

We compute back-of-the-envelope welfare loss estimates from cover % REVISION-V5: reframed as back-of-envelope per referee M5
bidding using the formula
$\text{Welfare loss} = (\exp(\hat{\beta}) - 1) \times
\sum_{i: \text{FL}_i = 1} p_i$.

We report four estimates using different coefficients. The OLS
coefficient (0.064) may be attenuated by measurement error in the binary
treatment indicator, but could also be inflated by positive selection
bias if FL firms participate in more expensive tenders within item
groups. The cross-fit coefficient (0.036) breaks any mechanical link
between FL classification and outcomes. The IV coefficient (0.194)
identifies a LATE for complier tenders, not an ATE.

Table~\ref{tab:welfare_bounds} in Appendix~\ref{app:extensions} reports
the results. The OLS-based estimate is R\$734 million (95\% CI:
R\$243M--R\$1,246M). The cross-fit estimate is R\$410 million. The
IV-based estimate is R\$2.4 billion, though the IV's fragility across
fixed-effects structures (Section~\ref{sec:results_iv}) warrants caution.
The concentrated-market FL coefficient ($-0.018$) produces a % REVISION-V5: discussed negative welfare per referee M5/minor 12
\emph{negative} welfare estimate (R\$$-88$M), consistent with the
finding that FL firms in concentrated markets show no price-inflating
effect---this subgroup does not contribute to the welfare cost of cover
bidding.

As a share of total procurement spending (R\$30.8 billion), the OLS
estimate represents approximately 2.4\% of total BEC procurement
value over 2009--2019 and the cross-fit estimate 1.3\%.
These figures exclude general equilibrium effects (deadweight loss,
allocative inefficiency, entry deterrence) and should be interpreted
as illustrative magnitudes rather than precise welfare bounds.

\subsection{Oversight Heterogeneity}

Table~\ref{tab:regime_oversight} in Appendix~\ref{app:robustness}
reports the FL price coefficient by PBU size quartile. The coefficient
is largest for the smallest PBUs (Q1: 0.214, imprecisely estimated) and
smallest for the largest PBUs (Q4: 0.017, not individually significant),
broadly consistent with the conceptual framework's prediction that
detection probability $\theta$ constrains cover bidding. Only Q3 (0.066)
is individually significant ($p < 0.05$). The pattern across quartiles
is noisy, reflecting limited variation in FL presence within small PBUs.

\begin{figure}[htbp]
\centering
\includegraphics[width=0.85\textwidth]{fig_07_threshold_stability.pdf}
\caption{FL price coefficient across IQR multiplier thresholds.
Error bars show 95\% confidence intervals.}
\label{fig:threshold_stability}
\end{figure}
